% Source: Paper_A_新版完整论文提纲_vFinal_Draft.md §18
\section{结论}

本文建立了一条从高信息 PreScreen 到前瞻路由验证的完整证据链。核心贡献是三轴工人测量、
双波 Calibration 与部分收缩、以及 Support-aware Full 对 Strong Global 的前瞻政策比较。
P1 形成可检验的 failure-family 假设；C1/C2 分离外部正确性、同行兼容性和结构可靠性，
并以 common anchor、diverse bridge 与精度补测冻结有限条件画像；T1 检验 Semi-Auto 在不同
模型风险下的质量与效率；V1 则在严格控制共享容量、offer、替补、动态冗余和聚合终态后，
比较支持度感知 Full 与 Strong Global。

论文的核心主张不是个体化必然优于全局策略，而是：

\begin{quote}
高信息诊断可以产生具有预测价值的条件假设；当这些假设在独立 Calibration 中得到支持并且任务侧可被可靠激活时，它们能够形成可审计的条件路由；当支持不足时，系统应明确退化为强 Global，而不是强行使用不稳定画像。
\end{quote}
